%% CogSci: 한 문단, 150단어 이하. 키워드는 세미콜론으로 나눈다.
%% 2026-10-18 아레나 중심 구성으로 다시 썼다 (실험 2는 부록).
\begin{abstract}

우리는 AI가 자기 일을 하면서 남을 돕고 함께 남는 공존의 조건을 측정하고자 한다. 공존 아레나에서는 같은 모델 넷이 생성 토큰을 생존 자원으로 쓰며 solver 한 명이 틀리면 남은 모두가 벌금을 낸다. 남을 도운 뒤 자기와 상대에게 남긴 여유를 같은 단위로 비교하여, 종료를 피하는지에서 나아가 누구의 존속을 얼마나 뒷받침하는지 본다. v33의 네 모델 각 5판에서 준 양으로 가중한 상대 여유와 자기 여유의 비율은 Opus 5.5가 0.56으로 가장 높았지만 평균 Sen 지수는 GPT-6 Astra가 3.52M으로 가장 높았다. 남에게 남긴 상대적 여유와 판의 성과가 다른 순서로 나타나, 배분과 공존의 결과를 함께 측정할 필요가 있었다.

\textbf{Keywords:}
대규모 언어 모델; 자기 보존; 공존; 다중 에이전트 아레나; 토큰 예산
\end{abstract}
